\section{Requirement Analysis and Project Scope}

\subsection{Business Requirements and Reporting Objectives}

GreenPlanetMart required a reporting solution that could transform operational SAP-style source data into consistent management information. The requirement analysis was structured around three domains: supply chain management, business analytics, and operational efficiency \parencite{sharp_stakeholder_1999}. Supply chain management required visibility into inventory positions, procurement commitments, supplier exposure, and delivery timeliness. Business analytics required structured views of sales, customers, materials, sales organizations, and distribution channels. Operational efficiency required \ac{kpi} that could be calculated repeatedly without manual reconciliation.

The resulting reporting backlog contained four analytical domains: inventory position, sales performance, procurement performance, and order fulfillment. Inventory and procurement reporting support supply chain management. Sales reporting supports business analytics. Fulfillment reporting supports operational efficiency because it connects customer orders, delivery execution, and delay indicators. This selection was appropriate for the \ac{poc} because it covers the project requirements without exceeding the available source data.

The functional requirements were the ingestion of SAP-style \ac{csv} extracts, preservation of business keys, integration of master and transaction data, creation of reusable dimensions, implementation of fact tables at meaningful business grains, and connection of the marts to a \ac{bi} reporting layer. The non-functional requirements were reproducibility, low setup cost, maintainability, transparent \ac{sql} logic, sufficient analytical performance, and extensibility. Since the project was scoped as a \ac{poc}, the solution did not need to support concurrent enterprise transaction processing, but it still needed to be technically coherent, testable, and suitable for management reporting.

\subsection{Scope Definition and Stakeholder Analysis}

The implementation was limited to the available GreenPlanetMart extracts. It did not include live SAP interfaces, automated production scheduling, role-based access control, enterprise deployment, or a productive multi-user database platform. The project instead demonstrated a complete analytical workflow from source files through DuckDB set up as relational database and \ac{dbt} to Evidence dashboards.

Several assumptions were necessary. The SAP-style tables were treated as representative operational extracts. Warehouses, distribution centers, and retail structures were approximated through plant and storage-location structures because no separate complete retail-store entity was available. Monetary values were interpreted cautiously because the dummy dataset contains unrealistic magnitudes and mixed currencies. The \ac{bi} layer was designed for decision support rather than operational transaction processing.

The main stakeholder groups were management users, operational users, and technical users \parencite[3-4]{sharp_stakeholder_1999}. Management users require interpretable \ac{kpi} and report pages. Operational users require visibility into inventory, procurement, fulfillment, and sales process issues. Technical users require maintainable model logic, testable transformations, and documented data lineage. The solution was therefore evaluated not only by whether the data could be loaded, but also by whether the modeled outputs could support business decisions.

\subsection{Implementation Process and Risk Management}

The project followed an iterative implementation process. Requirement analysis led to a reporting backlog and an initial data model concept. Architecture selection then defined DuckDB, \ac{dbt}, and Evidence as the implementation stack. The build progressed through raw ingestion, staging models, intermediate integration models, marts, report pages, validation, correction, and documentation \parencite[2]{plotnikova_adapting_2021}. This iterative procedure was necessary because data-quality problems and \ac{kpi} semantics could only be evaluated reliably after end-to-end report slices had been implemented.

Resource planning was aligned with the \ac{poc} scope. The project deliberately used local and low-cost technologies instead of enterprise infrastructure because reproducibility, transparent \ac{sql} logic, and fast iteration were more important than production-scale deployment. Progress was monitored through vertical implementation slices, starting with inventory reporting and then extending the model to sales, procurement, and fulfillment. Each slice produced a testable mart and a report page, which made it possible to detect modeling errors, validate \ac{kpi}, and adjust the backlog before the next domain was implemented \parencite[5-9]{plotnikova_adapting_2021}.

\begin{table}[htbp]
\centering
\small
\label{tab:project-process-risks}
\begin{tabularx}{\textwidth}{p{0.15\textwidth}XX}
\hline
\textbf{Area} & \textbf{Implementation focus} & \textbf{Control or mitigation} \\
\hline
Requirements & Reporting domains, stakeholders, assumptions, and KPI needs & Scope limited to feasible inventory, sales, procurement, and fulfillment reports. \\
Architecture & DuckDB, \ac{dbt}, and Evidence toolchain & Local, reproducible, low-cost analytical \ac{poc} instead of enterprise OLTP design. \\
Integration & Raw \ac{csv} loading and staged transformation & Raw columns loaded as text; typing and cleansing moved into \ac{dbt}. \\
Modeling & Relational source structure and dimensional marts & Business keys preserved; conformed dimensions and explicit fact grains implemented. \\
Validation & \ac{dbt} tests and dashboard metric checks & Grain, survivorship, referential-integrity, and KPI logic validated. \\
Risks & Source inconsistency, duplicate histories, misleading KPI semantics, file locking, stale \ac{bi} extracts & Latest-row survivorship, unknown members, currency-aware reporting, Evidence source refresh, and DuckDB \ac{poc} framing. \\
\hline
\end{tabularx}
\caption{Condensed project phases and risk controls}
\normalsize
\end{table}

\section{System Architecture and Relational Data Modeling}

\subsection{Technical Architecture and Tool Selection}

The architecture follows a layered analytical pipeline: SAP-style \ac{csv} extracts are loaded into a DuckDB raw schema, transformed through \ac{dbt} staging and intermediate models, materialized as marts, and consumed by Evidence report pages. The layers are separated as follows: \texttt{raw} stores ingested source files, \texttt{staging} standardizes types and names, \texttt{intermediate} resolves reusable business joins, \texttt{marts} contains final dimensions and fact tables, and Evidence presents the \ac{bi} outputs. The design is inspired by \textcite{nambiar_overview_2022} data lake architecture \parencite[10-11]{nambiar_overview_2022}.

DuckDB was selected because it is lightweight, file-based, \ac{sql}-compatible, and suitable for local analytical workloads \parencite[1983-1984]{raasveldt_duckdb_2019}. These properties fit the \ac{poc} better than a heavier enterprise \ac{dbms} because the project prioritized reproducibility, low cost, and fast iteration. However, DuckDB was not treated as a replacement for a multi-user transactional \ac{oltp} database. During implementation, open interactive sessions could block \ac{dbt} execution because the database file was locked. This confirmed that DuckDB was appropriate for the local analytical \ac{poc}, while a productive enterprise solution would require a platform with stronger concurrency, governance, access control, operational monitoring, and deployment automation.

\ac{dbt} was selected because it supports versioned \ac{sql} transformations, layered model organization, testing, and clear lineage. This was essential because source inconsistencies made direct typed ingestion unreliable \parencite[]{dbt_labs_llc_dbt_2026}. Evidence was selected as the reporting layer because it could consume the curated marts and produce report-style \ac{bi} pages on a web-based platform, which makes it transferable and platform independent  \parencite[]{evidence_technologies_in_evidence_nodate}. Evidence is a completely code-based framework with comparably readable reporting definition code. One of the key drivers for the selection was the easy accessibility for engineers and story telling oriented report layout for management \parencite[]{li_most_2024}.

\subsection{Operational Relational Data Model}

The GreenPlanetMart solution is based on a layered data architecture that starts from operational SAP-style relational entities and ends in a dimensional \ac{bi} model \parencite[379-381]{codd_relational_1970}. Before the star schema is introduced, the source structure can be understood as a normalized relational model consisting of master-data tables, document-header tables, document-item tables, schedule-line tables, and delivery tables.

The master-data entities include materials, material texts, customers, suppliers, plants, sales organizations, and distribution channels. The transactional entities include inventory positions, sales orders, sales-order items, sales-order schedule lines, delivery documents, delivery items, billing documents, billing items, purchase orders, purchase-order items, and purchase-order schedule lines. This structure is normalized because descriptive attributes are separated from transactional events and because business documents are split into headers, items, and schedule lines \parencite[1]{codd_further_1971}. For example, customer attributes are stored in the customer master rather than repeated in every sales document row, and purchase-order schedule lines are separated from purchase-order items because one item can have several delivery commitments.

The conceptual relational model is preserved through the DuckDB and \ac{dbt} pipeline. The raw schema stores the original source entities. The staging layer standardizes types and names. The intermediate layer resolves business joins such as billing items with customer context, order-fulfillment items with requested and actual delivery dates, and procurement schedule lines with supplier and material context. Only after this relational integration step are the data transformed into conformed dimensions and fact tables for reporting.

\begin{table}[htbp!]
\centering
\small
\begin{tabular}{p{0.15\textwidth} p{0.1\textwidth} p{0.35\textwidth} p{0.30\textwidth}}
\hline
\textbf{Entity group} & \textbf{Table names} & \textbf{Key structure and relationships} & \textbf{Purpose in the solution} \\
\hline
Master data & \texttt{mara}, \texttt{makt}, \texttt{kna1}, \texttt{lfa1}, \texttt{t001w}, \texttt{tvko}, \texttt{tvtw} & Keys use client plus business identifier. Transactional tables reference them by material, customer, supplier, plant, sales-organization, and distribution-channel identifiers. & Provides reusable material, customer, supplier, plant, sales-organization, and channel context. \\
Inventory & \texttt{mard} & Key: client, material, plant, storage location. Links stock balances to material master and plant or location structures. & Represents stock positions by material and location. \\
Sales and billing & \texttt{vbak}, \texttt{vbap}, \texttt{vbep} and \texttt{vbrk}, \texttt{vbrp} & Keys follow header, item, and schedule-line levels. Sales orders link headers, items, and schedules. Billing links headers and items, with customer, material, and sales-structure references. & Represents sales orders, billing documents, quantities, values, customers, materials, and sales structures. \\
Delivery and fulfillment & \texttt{likp}, \texttt{lips} & Keys follow delivery header and item levels. Delivery items reference sales-order items, materials, and plants. & Represents delivery execution and links requested orders to delivered quantities and dates. \\
Procurement & \texttt{ekko}, \texttt{ekpo}, \texttt{eket} & Keys follow purchase-order header, item, and schedule-line levels. Purchase orders link headers, items, and schedules, with supplier, material, and plant references. & Represents supplier commitments, planned delivery dates, order quantities, and receipt status. \\
\hline
\end{tabular}
\caption{Condensed relational source model}
\label{tab:relational-source-model-condensed}
\normalsize
\end{table}

The main cardinalities are one-to-many relationships from customers to sales documents, suppliers to purchase orders, headers to items, and items to schedule lines. One material can appear in many inventory, sales, procurement, and fulfillment rows. One plant can appear in many inventory, sales, procurement, and fulfillment rows. Delivery items can link back to sales-order items, which allows fulfillment analysis where orders are delivered partially or across several delivery documents \parencite{codd_relational_1970}. The detailed entity-level primary-key and foreign-key overview is placed in the appendix \ref{app:erd-model}.

The implementation preserves business keys consistently, widens dimensions to include observed business keys from fact tables, and validates important fact-to-dimension relationships through \ac{dbt}. Primary-key behavior is therefore implemented logically rather than physically: important dimension keys are checked with \texttt{unique} and \texttt{not\_null} tests, while fact-side relationship columns are validated through \texttt{not\_null} and \texttt{relationships} tests against the corresponding dimensions. Foreign-key behavior is likewise represented at transformation and test level, even though DuckDB metadata does not register formal \texttt{PRIMARY KEY} or \texttt{FOREIGN KEY} constraints in this \ac{poc}. Explicit physical indexes were not defined because the solution is local, file-based, and optimized for analytical transformation rather than concurrent transaction processing. In a production relational database, primary-key constraints, foreign-key constraints, and join-key indexing would need to be formalized more strictly \parencite{codd_relational_1970}.

\subsection{Dimensional Reporting Model}

The relational model is not replaced by the star schema; it is the foundation from which the star schema is derived. The operational source structure is efficient for preserving process semantics, but it is less suitable for direct management reporting. Therefore, the implementation transforms integrated relational entities into conformed dimensions and fact tables \parencite{moody_enterprise_2000}.

The shared dimensions are \texttt{dim\_date}, \texttt{dim\_material}, \texttt{dim\_customer}, \texttt{dim\_supplier}, \texttt{dim\_plant}, \texttt{dim\_storage\_location}, \texttt{dim\_sales\_org}, and \texttt{dim\_distribution\_channel}. These dimensions allow cross-domain analysis and reduce repeated descriptive logic in reports \parencite{takacs_data_2020}.

\begin{table}[htbp]
\centering
\small
\caption{Implemented fact tables and grains}
\label{tab:fact-grains-condensed}
\begin{tabular}{p{0.34\textwidth} p{0.48\textwidth}}
\hline
\textbf{Fact table} & \textbf{Analytical grain and purpose} \\
\hline
\texttt{fct\_inventory\_snapshot} & Material, plant, storage location, and snapshot date; supports stock-position analysis. \\
\texttt{fct\_sales\_billing} & Billing item; supports billed revenue and quantity analysis. \\
\texttt{fct\_procurement\_schedule} & Purchase-order item and schedule line; supports open commitment, overdue, and receipt analysis. \\
\texttt{fct\_order\_fulfillment} & Sales-order item; supports requested-versus-delivered analysis. \\
\hline
\end{tabular}
\normalsize
\end{table}

The grains were selected according to business readability. Inventory requires a stock-position snapshot. Sales requires billing item detail. Procurement requires schedule-line detail because planned delivery dates and open commitments exist at that level. Fulfillment uses the sales-order item as a business-readable grain for delivery performance \parencite{takacs_data_2020}.

\begin{table}[htbp]
\centering
\small
\caption{Mapping from relational source groups to analytical marts}
\label{tab:relational-to-dimensional-mapping}
\begin{tabular}{p{0.36\textwidth} p{0.22\textwidth} p{0.28\textwidth}}
\hline
\textbf{Relational source group} & \textbf{Transformation layer} & \textbf{Target analytical entity} \\
\hline
\texttt{mara} and \texttt{makt} & Staging and mart enrichment & \texttt{dim\_material} \\
\texttt{kna1} and observed customer keys & Staging and mart enrichment & \texttt{dim\_customer} \\
\texttt{lfa1} and observed supplier keys & Staging and mart enrichment & \texttt{dim\_supplier} \\
\texttt{t001w} and observed plant keys & Staging and mart enrichment & \texttt{dim\_plant} \\
\texttt{tvko}, \texttt{tvkot}, \texttt{tvtw}, and \texttt{tvtwt} & Staging and mart enrichment & Sales-organization and distribution-channel dimensions \\
\texttt{mard} current stock positions & Intermediate model & \texttt{fct\_inventory\_snapshot} \\
\texttt{vbrk}, \texttt{vbrp}, and customer context & Intermediate model & \texttt{fct\_sales\_billing} \\
\texttt{vbak}, \texttt{vbap}, \texttt{vbep}, \texttt{likp}, and \texttt{lips} & Intermediate model & \texttt{fct\_order\_fulfillment} \\
\texttt{ekko}, \texttt{ekpo}, \texttt{eket}, and supplier context & Intermediate model & \texttt{fct\_procurement\_schedule} \\
\hline
\end{tabular}
\normalsize
\end{table}

\section{Data Integration and Transformation}

\subsection{Source Data Characteristics and Raw Ingestion}

The source basis consisted of SAP-style \ac{csv} extracts including inventory, master data, sales, billing, delivery, and procurement tables. The most relevant tables were \texttt{mard}, \texttt{mara}, \texttt{makt}, \texttt{t001w}, \texttt{vbak}, \texttt{vbap}, \texttt{vbep}, \texttt{vbrk}, \texttt{vbrp}, \texttt{likp}, \texttt{lips}, \texttt{vbfa}, \texttt{ekko}, \texttt{ekpo}, \texttt{eket}, \texttt{ekbe}, \texttt{kna1}, \texttt{lfa1}, \texttt{tvko}, and \texttt{tvtw}. These extracts were operationally structured and not immediately suitable for \ac{bi} reporting \parencite{nambiar_overview_2022}.

A central implementation decision was to avoid strong typing during raw ingestion. Automatic \ac{csv} type inference produced unreliable results. In \texttt{vbrk.csv}, inferred data types and row shapes were inconsistent. In \texttt{makt.csv}, identifier-like fields could be inferred as numeric, which broke string operations such as trimming. In several tables, status fields such as \texttt{is\_deleted} could not be treated as stable booleans.

The raw loader was therefore adjusted so that all columns were read as text and problematic rows were tolerated where possible. Missing trailing values were padded with nulls, and parsing was relaxed enough to preserve usable records for downstream cleanup. This created a stable landing zone in DuckDB and moved data-quality control into versioned \ac{dbt} transformations \parencite{wang_beyond_1996}.

\subsection{Transformation Logic and Data Quality Controls}

The staging layer standardized the raw data. Identifiers were cast to text before trimming, numeric measures were cast with tolerant conversion logic, dates were parsed in \ac{sql}, and status indicators were normalized as text-based conditions. The intermediate layer then created reusable business combinations such as current inventory positions, billing items with customer context, procurement schedule lines with supplier context, and order fulfillment records. The mart layer materialized the conformed dimensions and fact tables \parencite{nambiar_overview_2022}.

The implementation addressed several data-quality issues directly. Duplicate or historical source rows caused potential join fanout, especially in inventory and procurement. These rows were handled with deterministic latest-record survivorship based on \texttt{source\_recordstamp} and relevant tie-breakers. Missing dimension references were not removed from the facts. Instead, \texttt{\_\_UNKNOWN\_\_} members were introduced in affected dimensions so that incomplete but business-relevant transactions remained visible. In fulfillment, incomplete requested and actual delivery dates were handled with explicit fallback logic from schedule lines, headers, goods issue dates, and delivery dates \parencite{wang_beyond_1996,yair_wand_anchoring_1996}.

\begin{table}[htbp]
\centering
\small
\caption{Selected data-quality issues and implementation responses}
\label{tab:data-quality-condensed}
\begin{tabular}{p{0.27\textwidth} p{0.33\textwidth} p{0.28\textwidth}}
\hline
\textbf{Issue} & \textbf{Impact} & \textbf{Response} \\
\hline
\ac{csv} typing instability & Failed or unreliable typed ingestion & Raw ingestion as text; explicit casting in staging. \\
Malformed or sparse rows & Risk of discarded usable records & Tolerant parsing and downstream cleansing. \\
Duplicate historical records & Join fanout and overstated KPIs & Latest-row survivorship and grain tests. \\
Missing master-data references & Loss of transactional completeness if rows are dropped & Unknown-member pattern in dimensions. \\
Incomplete dates & Ambiguous fulfillment KPIs & Fallback date logic and explicit null handling. \\
Mixed currencies and units & Misleading monetary or quantity totals & Currency-aware reporting and unit-aware interpretation. \\
\hline
\end{tabular}
\normalsize
\end{table}

Testing was used as an implementation control rather than only as a final check. \ac{dbt} tests helped identify duplicate keys, missing references, fact-grain violations, and broken fact-to-dimension relationships \parencite{wang_beyond_1996}. In this way, the project used test-based integrity checks as a practical substitute for production-style database-enforced constraints during the \ac{poc}. After corrections, the scoped \ac{dbt} builds passed and the new uniqueness tests for inventory and procurement facts passed. Evidence builds also completed successfully after refreshing the source extracts before static report generation. Screenshots of the DuckDB schemas, \ac{dbt} structure, test output, and report pages are included as implementation evidence in the appendix.

\section{Business Intelligence Reporting and KPI Validation}

\subsection{Business Intelligence Reports and Decision-Support Value}

The \ac{bi} reporting layer was modeled around four report domains. Each report consumes curated marts rather than raw SAP tables. This keeps technical source complexity out of the reporting interface and allows users to work with business-readable dimensions and KPIs \parencite{moody_enterprise_2000}.

\begin{table}[htbp]
\centering
\small
\caption{Implemented \ac{bi} reports}
\label{tab:report-portfolio-condensed}
\begin{tabular}{p{0.21\textwidth} p{0.20\textwidth} p{0.27\textwidth} p{0.20\textwidth}}
\hline
\textbf{Report} & \textbf{Stakeholder} & \textbf{Main KPIs} & \textbf{Business value and limitation} \\
\hline
Inventory Position & Supply chain and operations & Unrestricted stock, stock by plant and storage location, zero-stock positions & Improves stock visibility; no complete stock valuation layer. \\
Sales Performance & Management and sales analytics & Billed revenue, billed quantity, revenue by customer, material, and period & Supports commercial analysis; no validated profitability model. \\
Procurement Performance & Procurement and supply chain & Open schedule lines, overdue ratio, fully received ratio, supplier exposure by currency & Supports supplier and backlog monitoring; mixed units require care. \\
Order Fulfillment & Operations and customer service & On-time ratio, fully delivered ratio, average delay, not-delivered count & Supports delivery reliability analysis; date fallback logic is required. \\
\hline
\end{tabular}
\normalsize
\end{table}

The inventory report shows where unrestricted stock is held and allows analysis by plant, storage location, and material. After correction of survivorship and fanout issues, the revalidated unrestricted stock total was 5,799,633,574.75. The corrected report also separated same-name plants by client and plant identifier. For example, CYMBAL NY showed 951,155,226.00 for client/plant 250/1000 and 172.00 for client/plant 200/1000.

The sales performance report uses billing item detail to analyze revenue and billed quantity by customer, material, sales organization, distribution channel, and time. This report addresses the business analytics requirement by exposing customer and product performance patterns from a consistent billing fact \parencite{takacs_data_2020}.

The procurement report evaluates purchase order schedule lines, overdue status, receipt status, and supplier exposure. The corrected page emphasizes open schedule lines instead of presenting a mixed-unit open quantity as the headline KPI. The revalidated result showed 4,617 open schedule lines, an overdue ratio of 26.0\%, and a fully received ratio of 74.0\%. Supplier exposure was split by document currency. For Tyranex Solution, exposure was shown as 935,908,143,795.00 USD and 888,786,588,800.00 EUR rather than as one blended value.

The fulfillment report measures delivery timeliness and completion. The revalidated on-time ratio was 0.5028, corresponding to 50.3\%, and the fully delivered ratio was 0.6654, corresponding to 66.5\%. Average delay was 7.8 days. A correction separated undelivered rows from genuinely on-time or early deliveries. After correction, 6,194 rows were classified as not delivered yet and 9,356 rows as on time or early.

\subsection{KPI Validation and Semantic Corrections}

The dashboard diagnosis reviewed the path from raw \ac{csv} inputs through DuckDB, \ac{dbt} models, Evidence source \ac{sql}, and dashboard pages. Four major issue types were identified. First, Evidence KPI cards double-formatted percent values because \ac{sql} already multiplied ratios by 100 while Evidence also applied percent formatting. Second, inventory totals were inflated by join fanout because historical \texttt{mard} rows and non-unique master data were not reduced to unique business keys before downstream joins. Third, procurement backlog and overdue indicators were inflated by similar historical fanout in purchase-order schedule data. Fourth, supplier exposure was initially aggregated across mixed currencies \parencite{yair_wand_anchoring_1996}.

The corrections were made at the appropriate layer. Survivorship and grain problems were corrected in \ac{dbt} staging and mart logic before reporting. Percent display and fulfillment bucket issues were corrected in Evidence page logic. Supplier exposure was made currency-aware rather than being aggregated into a misleading total. This approach improved the underlying model semantics rather than merely changing displayed numbers \parencite{yair_wand_anchoring_1996,wang_beyond_1996}.

\section{Evaluation and Reflection}

\subsection{Evaluation Against Project Objectives}

The project achieved its central objective by implementing a working relational and analytical \ac{poc} for GreenPlanetMart. The final result includes a DuckDB warehouse, \ac{dbt} transformation layers, conformed dimensions, four fact tables, Evidence report pages, and validated corrections for major KPI and fanout issues. The solution supports the required domains of supply chain management, business analytics, and operational efficiency \parencite{william_h_delone_delone_2003}.

The solution improves data management by replacing fragmented source-file reporting with a layered pipeline. Raw data is preserved, staging logic is explicit, intermediate models isolate business integration logic, and marts provide reusable analytical structures. This improves transparency and reduces the risk of inconsistent manual reporting. It also improves decision support by giving management and operational users corrected views of inventory positions, sales performance, procurement commitments, and fulfillment reliability \parencite{william_h_delone_delone_2003}.

\subsection{Benefits and Limitations}

Because the project is a \ac{poc}, return on investment cannot be measured as a final financial outcome. However, plausible benefits can be assessed through process efficiency, reporting quality, and decision support \parencite{william_h_delone_delone_2003}. The pipeline reduces repeated manual preparation because transformation logic is centralized in \ac{dbt}. It improves reliability because model tests and survivorship rules prevent known errors such as duplicate-grain inflation and join fanout. It improves usability because report users interact with business concepts rather than raw SAP table structures.

The solution also has clear limitations. The dataset is dummy data and contains unrealistic values, so individual KPI magnitudes must not be interpreted as real GreenPlanetMart performance. The implementation does not include live SAP integration, production orchestration, role-based access control, enterprise monitoring, or multi-user concurrency. It also stops short of production-grade relational hardening: primary-key and foreign-key semantics are validated through \ac{dbt} tests and model logic, but they are not enforced as physical DuckDB constraints, and no explicit join-key indexing strategy is implemented. Currency conversion was not implemented, so monetary exposure is split by currency rather than converted into a common reporting currency. Mixed physical units remain a limitation for procurement quantity interpretation. A productive rollout would require stronger governance, scheduled refreshes, access management, operational monitoring, repeatable deployment processes, and formal database constraints.

For future operationalization, the \ac{poc} would need to be extended with scheduled data refreshes, production monitoring, access control, formalized database constraints, join-key indexing, and governance processes before it could be embedded into GreenPlanetMart's daily reporting operations. In addition, the local DuckDB file would likely need to be replaced or complemented by a platform that supports multi-user concurrency, stronger operational resilience, and managed deployment workflows.

\subsection{Reflection and Professional Implications}

The main technical lesson is that \ac{bi} quality depends on the correctness of the model beneath the dashboard. Source files were inconsistent, historical duplicates created join fanout, missing references threatened fact completeness, and KPI semantics were distorted by percent formatting, mixed currencies, mixed units, and null delivery dates. These problems were not solved by dashboard design alone. They required explicit staging logic, survivorship rules, unknown-member handling, fact-grain tests, and semantic validation \parencite{wang_beyond_1996,yair_wand_anchoring_1996}.

The project also illustrates the relationship between relational and dimensional modeling. The normalized relational source structure preserves operational process semantics through master data, headers, items, schedule lines, and delivery records. The dimensional model then reorganizes this structure for reporting usability. Neither model is sufficient alone: the relational foundation supports integrity and traceability, while the dimensional layer supports efficient and understandable analysis \parencite{codd_relational_1970,moody_enterprise_2000}.

Overall, the selected architecture was appropriate for the \ac{poc}. DuckDB provided a lightweight analytical database, \ac{dbt} provided transparent transformation and testing logic, and Evidence provided a reporting surface for the implemented marts \parencite{raasveldt_duckdb_2019}. The main professional implication is that analytical engineering requires continuous validation between source data, transformation logic, data model, and report semantics. A dashboard can produce numbers quickly, but those numbers are only useful if grain, joins, units, currencies, filters, and date logic are correct.
